\documentclass[12pt]{article}
\usepackage{notes-common}
\usetikzlibrary{calc}
\geometry{margin=0.35in}
\AtBeginDocument{\fontsize{14}{17}\selectfont}

% Complete / filled-in reference version of biot-savart-intro-notes.tex.
% Fri Oct 23 -- Magnetic fields + Biot-Savart (Knight 29.1-29.3).
% Built from Seth's rough notes (uploaded).
\newcommand{\ANS}[1]{\textcolor{accent}{#1}}
\newcommand{\disc}[1]{{\small\color{accent} #1}}

% A compact filled-in "(a, b, c)" vector triple, same visual slot as \vecbl.
% \textcolor works fine inside a single continuous math-mode span; the bug to
% avoid is closing and reopening $...$ around each component (that drops back
% to text mode between them, where \times etc. aren't valid).
\newcommand{\vecans}[3]{$\big(\textcolor{accent}{#1}\,,\ \textcolor{accent}{#2}\,,\ \textcolor{accent}{#3}\,\big)$}

% ---- same axis + point-charge diagram pieces as the skeleton ----
\newcommand{\axesxy}[1][2.6]{%
  \draw[-latex, line width=0.6pt, gray] (-#1,0) -- (#1,0) node[below right,scale=0.7]{$x$};
  \draw[-latex, line width=0.6pt, gray] (0,-#1) -- (0,#1) node[above left,scale=0.7]{$y$};
}
\newcommand{\qdot}[2]{
  \filldraw[fill=accent!25, draw=accent] (#1) circle (2.6pt);
  \node[scale=0.75, anchor=north, yshift=-2pt] at (#1) {#2};
}
\newcommand{\pdot}[2]{
  \filldraw[fill=white, draw=black] (#1) circle (2pt);
  \node[scale=0.75, anchor=south, yshift=2pt] at (#1) {#2};
}
\newcommand{\sepvec}[2]{\draw[-latex, accent, line width=0.8pt] (#1) -- (#2);}

\begin{document}

\noteshead{Intro to Biot--Savart \textcolor{accent}{(completed)}}{}

\goal{
\begin{itemize}[itemsep=5pt,parsep=2pt,topsep=2pt]
\item Warm up with something you already know: superpose the $\vec E$ field
of a point charge at three locations, being explicit about separation
vectors and unit vectors.
\item Reuse that exact same geometry to introduce the Biot--Savart law for a
single \emph{moving} point charge, $\vec B$.
\item Compare $\vec E$ and $\vec B$ at those same three points --- same
source, same locations, very different field.
\item Scale up to three moving charges in a line and see their $\vec B$
contributions add (not cancel) at an off-axis point.
\item Make the leap from one moving charge ($q\vec v$) to a continuous
current ($I\,d\vec\ell$) --- the form of Biot--Savart we'll actually use for
wires, loops and solenoids.
\end{itemize}}

%============================================================
\nsec{Warm-up: $\vec E$ from a point charge}

A charge $q=+2.0\ \mu$C sits at $(-5,0,0)$ cm. Find $\vec E$ at three field
points: $P_1=(-9,3,0)$, $P_2=(-5,3,0)$, $P_3=(-1,3,0)$ (all in cm).

\begin{center}
\begin{tikzpicture}[x=0.5cm,y=0.5cm]
  \axesxy[5.5]
  \sepvec{-5,0}{-9,3} \sepvec{-5,0}{-5,3} \sepvec{-5,0}{-1,3}
  \qdot{-5,0}{$q$}
  \pdot{-9,3}{$P_1$}
  \pdot{-5,3}{$P_2$}
  \pdot{-1,3}{$P_3$}
\end{tikzpicture}
\end{center}
\figtask{Sketch the three separation vectors $\vec r_1,\vec r_2,\vec r_3$ from
$q$ to each field point on the diagram above.}

\textbf{Step 1 --- separation vectors} $\vec r_i = \vec r_{P_i}-\vec r_q$ (cm):

$\vec r_1 = $ \vecans{-4}{3}{0}

$\vec r_2 = $ \vecans{0}{3}{0}

$\vec r_3 = $ \vecans{4}{3}{0}

\textbf{Step 2 --- magnitudes and unit vectors} $\hat r_i = \vec r_i / r_i$:

$r_1 = $ \ANS{5.0} cm, $\ \hat r_1 = $ \vecans{-0.8}{0.6}{0}

$r_2 = $ \ANS{3.0} cm, $\ \hat r_2 = $ \vecans{0}{1.0}{0}

$r_3 = $ \ANS{5.0} cm, $\ \hat r_3 = $ \vecans{0.8}{0.6}{0}

\textbf{Step 3 --- Coulomb's law}, $\displaystyle \vec E = \frac{1}{4\pi
\varepsilon_0}\,\frac{q}{r^2}\,\hat r$:

$\vec E_1 = $ \vecans{-5.75\times10^6}{4.32\times10^6}{0} N/C

$\vec E_2 = $ \vecans{0}{2.00\times10^7}{0} N/C

$\vec E_3 = $ \vecans{5.75\times10^6}{4.32\times10^6}{0} N/C

\checkpoint{$P_1$ and $P_3$ are the same distance from $q$. What's the same
about $\vec E_1$ and $\vec E_3$, and what's different? Why?}

\disc{Same magnitude ($7.19\times10^6$ N/C) and the same $+y$ component
($4.32\times10^6$ N/C), since $P_1$ and $P_3$ are the same height above $q$.
They differ only in the sign of the $x$-component --- $\vec E_1$ points up
and to the \emph{left}, $\vec E_3$ up and to the \emph{right}. That's exactly
mirror symmetry: $P_1$ and $P_3$ are mirror images of each other across the
vertical line through $q$.}

%============================================================
\nsec{Biot--Savart: the same charge, now moving}

Same charge, same three points --- but now $q$ moves with velocity
$\vec v = (12,0,0)$ m/s. The Biot--Savart law for a single point charge is

\[
\vec B = \frac{\mu_0}{4\pi}\,\frac{q\,\vec v\times\hat r}{r^2}, \qquad
\frac{\mu_0}{4\pi} = 1\times10^{-7}\ \text{T}\!\cdot\!\text{m/A}.
\]

\textbf{Reuse your $\hat r_1,\hat r_2,\hat r_3$ from \S1} --- the geometry
hasn't changed, only the source. Compute $\vec v\times\hat r_i$ first, then
$\vec B_i$:

$\vec v\times\hat r_1 = $ \vecans{0}{0}{7.2} m/s

$\vec v\times\hat r_2 = $ \vecans{0}{0}{12} m/s

$\vec v\times\hat r_3 = $ \vecans{0}{0}{7.2} m/s

$\vec B_1 = $ \vecans{0}{0}{5.76\times10^{-10}} T

$\vec B_2 = $ \vecans{0}{0}{2.67\times10^{-9}} T

$\vec B_3 = $ \vecans{0}{0}{5.76\times10^{-10}} T

\checkpoint{Is $\vec B_1$ the mirror image of $\vec B_3$, the way $\vec E_1$
and $\vec E_3$ were? What's actually true about $\vec B_1$ and $\vec B_3$?}

\disc{No --- $\vec B_1$ and $\vec B_3$ come out \emph{identical}, not mirror
images. $\vec v$ points along $\hat x$, so $\vec v\times\hat r$ only ever
depends on the $y$-component of $\hat r$ (the $x$-component of $\hat r$
drops out of that cross product entirely). $\hat r_1$ and $\hat r_3$ have
the \emph{same} $y$-component (0.6), so they give the same $\vec B$ even
though their $x$-components have opposite signs. $\vec B$ simply doesn't
see the left/right distinction that $\vec E$ did.}

%============================================================
\nsec{Compare $\vec E$ and $\vec B$}

\begin{itemize}[itemsep=6pt,parsep=2pt,topsep=2pt]
\item Which field is bigger here, $\vec E$ or $\vec B$ --- by roughly how
many orders of magnitude? \ANS{$E$ is bigger by about $10^{16}$. That's not
a coincidence: $\vec B = (\vec v\times\vec E)/c^2$, and $v/c^2 \approx
12/(9\times10^{16})\approx1.3\times10^{-16}$ --- exactly the ratio you'd
compute from $B_2/E_2$ above.}
\item $\vec E_1,\vec E_2,\vec E_3$ point in \ANS{the same directions as
$\hat r_1,\hat r_2,\hat r_3$ --- radially away from $q$, since $q$ is
positive.} directions (how do they relate to $\hat r$?).
\item $\vec B_1,\vec B_2,\vec B_3$ all point \ANS{straight out of the page,
$+\hat z$ --- because $\vec v$ ($\hat x$) and every $\hat r$ here (in the
$xy$-plane) both lie in the $xy$-plane, and the cross product of two
in-plane vectors is always perpendicular to that plane.} (what direction, and
why does every one of them end up that way)?
\end{itemize}

\checkpoint{$\vec v$ points along $+\hat x$ and every field point here has
$z=0$. Could $\vec B$ ever have come out with a nonzero $x$- or
$y$-component at these points? Why or why not?}

\disc{No. With $\vec v=(v,0,0)$ and $\hat r=(r_x,r_y,0)$ confined to the
$xy$-plane, $\vec v\times\hat r=(0,0,vr_y)$ --- the $x$- and $y$-components
of the cross product are identically zero no matter what $r_x$ and $r_y$
are. $\vec B$ \emph{has} to point along $\pm\hat z$ for any field point in
this plane; only its sign and size can change.}

%============================================================
\nsec{Three moving charges in a line}

Three charges, each $q=+1.0\ \mu$C, sit at $(1,-4,0)$, $(1,0,0)$, $(1,4,0)$
cm, each moving with $\vec v = (0,10,0)$ m/s. Find the total $\vec B$ at
field point $P=(-2,0,0)$ cm --- off to the side, not on the line of charges.

\begin{center}
\begin{tikzpicture}[x=0.45cm,y=0.45cm]
  \axesxy[6]
  \sepvec{1,-4}{-2,0} \sepvec{1,0}{-2,0} \sepvec{1,4}{-2,0}
  \qdot{1,-4}{$q_1$}
  \qdot{1,0}{$q_2$}
  \qdot{1,4}{$q_3$}
  \pdot{-2,0}{$P$}
\end{tikzpicture}
\end{center}
\figtask{Sketch the three separation vectors from each charge to $P$.}

$\vec r_1 = $ \vecans{-3}{4}{0} \quad $r_1 = $ \ANS{5.0} cm \quad
$\hat r_1 = $ \vecans{-0.6}{0.8}{0}

$\vec r_2 = $ \vecans{-3}{0}{0} \quad $r_2 = $ \ANS{3.0} cm \quad
$\hat r_2 = $ \vecans{-1.0}{0}{0}

$\vec r_3 = $ \vecans{-3}{-4}{0} \quad $r_3 = $ \ANS{5.0} cm \quad
$\hat r_3 = $ \vecans{-0.6}{-0.8}{0}

\textbf{Each charge's contribution} $\vec B_i = \dfrac{\mu_0}{4\pi}\dfrac{q\,
\vec v\times\hat r_i}{r_i^2}$:

$\vec B_1 = $ \vecans{0}{0}{2.4\times10^{-10}} T \qquad
$\vec B_2 = $ \vecans{0}{0}{1.11\times10^{-9}} T

$\vec B_3 = $ \vecans{0}{0}{2.4\times10^{-10}} T

\result[1.3cm]{$\vec B_{\text{total}} = \vec B_1+\vec B_2+\vec B_3 =
\ANS{(0,\,0,\,1.59\times10^{-9})\text{ T}}$}

\checkpoint{Did the three contributions partly cancel, or did they add? Compare
to what you'd expect from three point \emph{charges'} $\vec E$ fields at an
off-axis point --- is $\vec B$ behaving the same way?}

\disc{They added, all three pointing $+\hat z$. $\vec v$ here is $\hat y$,
so $\vec v\times\hat r=(0,0,-v\,r_x)$ --- it depends only on $\hat r$'s
$x$-component. All three charges sit to the \emph{right} of $P$, so all
three $\hat r_i$ have the same-sign $x$-component, and their contributions
add instead of partly canceling. Three point \emph{charges'} $\vec E$ fields
at this same off-axis $P$ would \emph{not} all add this cleanly --- $E_1$
and $E_3$ would have oppositely-signed $y$-components that partly cancel,
the same way $E_1$ and $E_3$ differed back in \S1. $\vec B$ is behaving
differently here because the cross product throws away a different
component than Coulomb's law keeps. This is exactly why a straight
current-carrying wire's field just builds up as you add more current
elements, instead of partly canceling, as you saw in the sim.}

%============================================================
\nsec{From one moving charge to a current}

A current is just a huge number of moving charges, confined to a wire. Take
a short piece of wire of length $d\ell$, carrying current $I$: in time $dt$,
the charge that passes through it is $dq$, moving at the drift velocity
$\vec v$ along the wire. The replacement rule is

\[
q\,\vec v \ \longrightarrow\ \ANS{I\,d\vec\ell}
\]

so the Biot--Savart law for one point charge, $\vec B = \dfrac{\mu_0}{4\pi}
\dfrac{q\,\vec v\times\hat r}{r^2}$, becomes the law for one current
element:

\result[2.4cm]{$d\vec B = \ANS{\dfrac{\mu_0}{4\pi}\,\dfrac{I\,d\vec\ell
\times\hat r}{r^2}}$}

\checkpoint{Everything you did in \S1--\S4 above with separation vectors and
unit vectors --- does any of it change once $q\vec v$ becomes $I\,d\vec
\ell$? What's actually new?}

\disc{None of the vector mechanics changes --- you still find a separation
vector, divide by its length to get $\hat r$, cross it with a direction
vector, and divide by $r^2$. What's new is bookkeeping, not mechanics:
instead of one discrete charge's $q\vec v$, a real wire has charge
distributed continuously along its length, so the \emph{sum} over point
charges becomes an \emph{integral} over the wire --- same ingredients, added
up continuously instead of one at a time. That's exactly the sum you built
by hand, element by element, in the simulator.}

\end{document}
